% Supplementary S16 — Filter-evaluation session: design and comparator
% Holds the over-specific 2nd-MRI design details moved out of Methods §Filter evaluation:
% macOS comparator specifics, run-count rationale (four per condition), single-case
% inference detail, and the behavioral battery.

\section*{Supplementary S16: Filter-evaluation session --- design and comparator}
\label{app:filter_eval}

\paragraph{Comparator.}
The deployed comparator was the macOS accessibility Color Filter (System Settings $>$ Accessibility $>$ Display, build 26.5.1), set at an intensity the participant self-tuned before scanning. It is the unmodified shipping product --- an operating-system-level post-display transform. We compared against this deployed filter rather than a re-implemented retinal-simulation transform so that the comparator is the correction an end user actually encounters. The personalized filter, by contrast, was rendered directly in PsychoPy from the frozen per-subject pre-image.

\paragraph{Acquisition.}
Eight fMRI runs were acquired in ABBA-counterbalanced order, split evenly into four runs per filter. Scanner time set this count: a within-session two-condition design at six runs per condition would run roughly 90 minutes, beyond the 60--75-minute window over which participant attention and head motion stay acceptable. Four runs per condition ($\sim$60 minutes) sits at the recommended boundary.

\paragraph{Run-count adequacy.}
A subsampling analysis on the six-run Session-1 data --- evaluating signal-quality, geometric-stability, and outcome-detection metrics across all $\binom{6}{n}$ run subsets for $n = 2$--$6$ --- indicated that four runs retains the single-subject landmark findings, with five preferred and four the acceptable fallback. We extended the same subsampling to the primary interpolation metric. Across all $\binom{6}{n}$ subsets, hV4 LOCO adjacent accuracy retained the Session-1 landmark at every run count down to four (Figure~\ref{fig:adjacc_saturation}): the healthy-control mean stayed above the $3/8$ chance level (at $n = 4$, $0.45$ versus $0.46$ at six runs), while both CVD participants stayed well below it (deutan $0.23$, protan $0.14$; single-case $d_{cc} < -2$ in both). The metric reported in the main text is therefore stable at the deployed run count. Four runs is, separately, below the count that simulations recommend for fine-grained RDM-based model discrimination, so the second session is powered to detect movement toward or away from the HC reference, not to adjudicate model classes. For every neural index, the Session-1 HC reference runs were subsampled to four to match the per-condition run count, so condition and HC estimates rest on equal data.

\begin{figure}[tb]
  \centering
  \includegraphics[width=\linewidth]{figS16_adjacc_saturation}
  \caption{\textbf{The primary interpolation metric is stable under run-count reduction.} LOCO adjacent accuracy (canonical FE-6 basis, OLS pseudoinverse decoder) as a function of the number of runs $n$, computed over all $\binom{6}{n}$ run subsets of the Session-1 data, per ROI. Black line, healthy-control mean $\pm$ SEM (shaded band); dashed line, the $3/8$ chance level; markers, the two CVD single cases (sub-08 deutan, sub-09 protan) and the deutan control sub-10. At hV4 the healthy-control mean stays above chance and the CVD deficit persists at every run count, including the deployed $n = 4$ (orange guide); the metric does not depend on having the full six runs.}
  \label{fig:adjacc_saturation}
\end{figure}

\paragraph{Single-case inference.}
Each condition's neural index was compared to the HC distribution as a standardized single-case effect size (Cohen's $d$ against the HC mean and SD), reported in place of an inferential $p$ because the grand-mean permutation null is biased for this design. Robustness was assessed in two ways: by leave-one-run separation (the contrast holding after any single run is dropped) and by re-extracting both conditions on the no-filter baseline voxel set.

\paragraph{Behavioral battery.}
JND staircases were run for eight hue pairs, and the eight-alternative forced-choice (8AFC) color-identification task comprised 64 trials. Both were compared to the unfiltered Session-1 baseline and to the HC distribution.
